% ID: SEL-D0023-OA-E01
% ESTADO: validada. S024, cotejada con el examen.
\textbf{a)} Las restricciones equivalen a
$(5x+8)/13\le y\le\min(x,20-x)$. La cota inferior corta a
$y=x$ en $(1,1)$ y a $y=20-x$ en $(14,6)$.
Las dos cotas superiores se cruzan en $(10,10)$. Entre $x=1$
y $x=14$ se obtiene el triángulo cerrado de la figura.
\begin{center}
\setlength{\unitlength}{4mm}
\begin{picture}(19,14)(-1,-1)
\qbezier(0,0)(8,0)(16,0)\put(16.2,0){$x$}
\qbezier(0,0)(0,5.5)(0,11)\put(0,11.2){$y$}
\qbezier(1,1)(5.5,5.5)(10,10)
\qbezier(10,10)(12,8)(14,6)
\qbezier(14,6)(7.5,3.5)(1,1)
\put(1.2,0.8){\makebox(0,0)[lt]{$(1,1)$}}
\put(10,10.5){\makebox(0,0)[b]{$(10,10)$}}
\put(14.2,6){\makebox(0,0)[l]{$(14,6)$}}
\end{picture}
\end{center}
\textbf{b)} En $(3,2.5)$ tenemos $x+y=5.5\le20$,
$x-y=0.5\ge0$ y $5x-13y+8=-9.5\le0$. Cumple las tres condiciones.

\textbf{c)} Como $x-y\ge0$, el mínimo es $\boxed{6}$ en todo
el segmento $(x,x)$ con $1\le x\le10$.
Una función lineal alcanza el máximo en algún vértice.
Los dos primeros dan 6 y $(14,6)$ da $\boxed{14}$.
Como este último es el único vértice con ese valor, es el único máximo.
